\textbf{Proof.}
Let \(\mathfrak G_{\mathcal P}\) be the product of the within-block permutation groups.  For \(A\in\mathcal C_n\), define its orbit average
\[
\overline A=|\mathfrak G_{\mathcal P}|^{-1}
\sum_{g\in\mathfrak G_{\mathcal P}}U_gAU_g^\top.
\tag{1}
\]
The covariance hull and every score-quality shell for \(S_{\mathcal P}\) are invariant under this action.  Since \(A\mapsto q_\rho(A;S_{\mathcal P})\) is a maximum of linear functionals, it is convex, and Jensen's inequality gives
\[
q_\rho(\overline A;S_{\mathcal P})\leq q_\rho(A;S_{\mathcal P}).
\tag{2}
\]

For an assignment \(z\), write \(r_j=\sum_{i\in G_j}z_i\).  Conditional on \(r\), uniform within-block permutation gives
\[
\mathbb E[z_i z_{i'}\mid r]=\frac{r_jr_\ell}{m_jm_\ell}
\qquad
(i\in G_j,\ i'\in G_\ell,\ j\ne\ell),
\tag{3}
\]
and, for distinct \(i,i'\in G_j\),
\[
\mathbb E[z_i z_{i'}\mid r]
=\frac{r_j^2-m_j}{m_j(m_j-1)}.
\tag{4}
\]
The second identity follows by expanding \(r_j^2\).  Equations (3)--(4) are exactly the entries of \(A_{\mathcal P}(rr^\top)\).  Thus a convex representation of \(A\) supplies \(G\in\mathcal G_{\mathcal P}\) with \(\overline A=A_{\mathcal P}(G)\), proving the full-hull converse by (2).

Conversely, express \(G\) as a mixture of \(rr^\top\) over the block-sum lattice.  For every such \(r\), choose signs with those block sums, permute them uniformly within blocks, and multiply the whole assignment by an independent fair sign.  This is a mean-zero law in \(\mathcal D_n\), and (3)--(4) show that its covariance is \(A_{\mathcal P}(rr^\top)\).  Mixing realizes \(A_{\mathcal P}(G)\).  Compactness proves attainment and hence the exact orbit reduction.

For parity, if every \(m_j\) is even, the lattice contains \(r=0\); its covariance annihilates all block-constant centered vectors, so \(\kappa(S_{\mathcal P})=0\).  Conversely, if a cut-hull mixture has zero compression on \(S_{\mathcal P}\), positivity implies that every supported assignment satisfies \(x^\top z=0\) for every \(x\in S_{\mathcal P}\).  Its block-sum vector \(r\) is therefore orthogonal to every \(t\) satisfying \(\sum_jm_jt_j=0\), hence is proportional to \((m_1,\ldots,m_B)\).  Balance gives \(\sum_jr_j=0\), so the proportionality constant is zero.  Thus every block sum is zero, which is possible only when every block size is even.

We now solve the unequal two-block problem.  Write the block sizes as \(m<\ell\).  Balance forces \(r=(R,-R)\), where
\[
R\in\{-m,-m+2,\ldots,m\}.
\]
Because \(m+\ell=n\) is even, the block sizes have the same parity.  The convex hull of the possible squares is therefore
\[
q=\mathbb E[R^2]\in[e,m^2],
\qquad
e=0\ \text{for even--even},
\qquad
e=1\ \text{for odd--odd}.
\tag{5}
\]
Every point of this interval is obtained by mixing squared lattice values.

Let \(w\) be the unit vector in the one-dimensional score space, constant on each block.  Direct normalization gives the score eigenvalue
\[
a(q)=w^\top A_{\mathcal P}(G)w=\frac{nq}{m\ell}.
\tag{6}
\]
On the within-block contrast space of a block of size \(j>1\), the eigenvalue is
\[
b_j(q)=\frac{j^2-q}{j(j-1)}.
\tag{7}
\]
These mutually orthogonal spaces decompose \(H_n\).  If \(m>1\), put \(B(q)=\max\{b_m(q),b_\ell(q)\}\).  For \(s=\rho^2\), allocating the nonscore energy to the larger contrast eigenvalue yields
\[
q_{\sqrt{s}}(A_{\mathcal P}(q);S_{\mathcal P})
=\max\{a(q),\,sa(q)+(1-s)B(q)\}.
\tag{8}
\]

The three eigenvalues meet at the complete-randomization moment
\[
q_c=\frac{m\ell}{n-1},
\qquad
a(q_c)=b_m(q_c)=b_\ell(q_c)=c_n.
\tag{9}
\]
For \(e\leq q\leq q_c\), subtraction of the affine expressions (6)--(7) gives \(b_m(q)\geq b_\ell(q)\geq a(q)\); for \(q\geq q_c\), \(a(q)\) is the largest eigenvalue.  Consequently the minimization of (8) reduces on \([e,q_c]\) to
\[
sa(q)+(1-s)b_m(q).
\tag{10}
\]
Its slope changes sign at
\[
s_0=\frac{\ell}{m(n-1)}.
\tag{11}
\]
Thus (10) is minimized at \(q_c\) for \(s\leq s_0\), and at \(q=e\) for \(s\geq s_0\).  Substituting those endpoints proves
\[
v^*(\sqrt{s};S_{\mathcal P})
=\min\{c_n,(1-s)b_0+s\kappa_2\},
\qquad
b_0=\frac{m^2-e}{m(m-1)},
\qquad
\kappa_2=\frac{ne}{m\ell}.
\tag{12}
\]
The branches meet at \(s_0\).  The endpoint laws in the statement realize \(q=e\), while complete randomization realizes \(q=q_c\).  Equation (12) and the admissibility-frontier theorem give \(\gamma(S_{\mathcal P})=s_0\) and \(\rho^*(S_{\mathcal P})=\sqrt{s_0}\).

We next solve regret exactly.  For fixed \(q\), (8) is a maximum of affine functions of \(s\), while the oracle value in (12) is concave in \(s\).  Therefore the regret is convex in \(s\), and its maximum on \([0,1]\) occurs at an endpoint.  Since \(v(0)=c_n\) and \(v(1)=\kappa_2\),
\[
\Delta_{\mathrm{reg}}(A_{\mathcal P}(q);S_{\mathcal P})
=\max\{\max(a(q),B(q))-c_n,\ a(q)-\kappa_2\}.
\tag{13}
\]
On \([e,q_c]\), the first term is \(b_m(q)-c_n\), which decreases, and the second is \(a(q)-\kappa_2\), which increases.  They are equal at
\[
q_{\mathrm{reg}}=(1-s_0)e+s_0q_c,
\tag{14}
\]
and their common value is
\[
s_0(c_n-\kappa_2).
\tag{15}
\]
Outside this interval, (13) is no smaller.  The orbit reduction makes (13)--(15) a converse for every covariance in \(\mathcal C_n\), not merely for invariant designs.  Mixing the parity-endpoint law with complete randomization using weight \(s_0\) gives (14), so the converse is attained.  Scaling (15) by \(4M/n\) proves the regret formula.

If \(m=1\), equation (5) forces \(q=1=q_c\).  Only the \(\ell\)-block contrast space remains, and (6)--(7) both equal \(c_n\).  Hence the loss is \(c_n\) at every quality, \(\kappa=c_n\), \(\gamma=1\), \(\rho^*=1\), and regret is zero.

Finally, for \(B\) equal even blocks, averaging also over block permutations reduces a covariance to \(aP_S+b(P_H-P_S)\), with \(pa+hb=n\).  The score-shell loss is \(\max\{a,sa+(1-s)b\}\).  When \(a\leq c_n\), one has \(b\geq a\) and
\[
sa+(1-s)b
=\frac{n}{h}(1-s)+a\left\{s-\frac{p}{h}(1-s)\right\}.
\]
The coefficient of \(a\) changes sign at \(s=p/(n-1)\), whereas \(a>c_n\) already has loss above the complete-randomization value.  Minimization therefore gives
\[
v^*(\sqrt{s};S_B)=\min\{c_n,\delta_m(1-s)\},
\qquad
\delta_m=\frac{n}{h},
\qquad
\beta=\frac{p}{n-1}.
\tag{16}
\]
Complete randomization and independent within-block complete randomization realize the two endpoints.  Mixing them with weight \(\beta\) on complete randomization has normalized regret \(c_n\beta=np/(n-1)^2\).  The endpoint inequalities \(r\geq L-c_n\) and \(r\geq K\), together with \(n\leq pK+hL\), give the same lower bound for every covariance.  This proves the homogeneous regret formula.  Setting \(B=n/2\) and block size two gives the displayed pair specialization. \(\square\)